\section{Discussion}
\label{sec:discussion}

The central finding of this paper is that uncertainty estimation alone is
not enough. Without a decision rule to act on the uncertainty signal, the
information it carries goes unused. The gap between ``having uncertainty''
and ``using it to reduce errors'' is the gap this work measures---and it
is large. TTA with adaptive deferral removes roughly 80\% of segmentation
errors; the same TTA uncertainty, absent any deferral policy, removes 0\%.
The predictions are identical; only the decision changes. Choosing TTA
over MC Dropout improves Unc-AUROC by roughly 16 points, but the deferral
policy determines how much of that improvement actually reaches the
clinician. Neither choice alone reaches the best operating point; both
matter, and they interact. Figure~\ref{fig:summary_scatter} makes this
visible: the six method--policy configurations span a wide region of the
design space, and the best position---upper-left, with high error
reduction at low deferral rate---requires getting both choices right.
This is why uncertainty is best treated as a property of a decision
process rather than of a model in isolation.

\paragraph{Why overconfident errors are the dangerous ones.} Calibration
metrics treat every miscalibrated bin equally. Deferral does not. In
retinal screening, the errors that matter are the confident
false negatives---thin vessels the model misses while assigning a
confident ``background'' label. These pixels have \emph{low} uncertainty
by construction, and no amount of recalibration changes that. A deferral
rule informed only by calibrated probability gives such a pixel a free
pass. The decision-aware view exposes the gap: calibration can be good
while safety is poor, because safety depends on \emph{which} errors are
confident, not on their average. Confidence-aware deferral partially
addresses this by deferring pixels whose predictions sit near the
decision boundary, but the fundamental limit---confident errors are
invisible to uncertainty-based rules---remains the single most important
failure mode to communicate to a deploying clinician.

\subsection{Why TTA produces better uncertainty than MC Dropout}

The large Unc-AUROC gap between TTA (0.881) and MC Dropout (0.722)
is the largest effect in our experiments.

MC Dropout approximates a posterior over weights by sampling binary dropout
masks. The resulting uncertainty reflects how much predictions change when
random subsets of neurons are disabled. This is a global perturbation: it
affects the entire network, and its effect on any given pixel depends on how
that pixel's representation propagates through all layers. The result,
empirically, is spatially smooth uncertainty. Entire regions receive similar
uncertainty values because nearby pixels share most of their computational
path through the encoder.

TTA, by contrast, perturbs the input. A horizontal flip changes which pixels
fall on vessel boundaries; a 90-degree rotation changes the orientation of
thin vessels relative to convolutional filters. The resulting disagreement is
spatially specific: it concentrates on pixels whose predictions depend on
local geometric context---exactly where segmentation errors occur.

A subtler factor: MC Dropout uncertainty is computed as mutual information
(Equation~\ref{eq:mi}), which requires per-pass predictions to disagree.
But dropout perturbations are often too small to flip a prediction: if the
model is confident that a pixel is background ($\hat{p} \approx 0.05$),
30 dropout passes might produce values in $[0.03, 0.08]$, all on the same
side of the decision boundary. The mutual information is low even though the
pixel might be wrong. TTA augmentations, being geometric transformations,
can produce larger disagreements because they fundamentally change which
input features align with which filter orientations.

Effective sample count also matters. With $T = 30$ MC Dropout passes, the
effective number of independent samples is lower than 30 because dropout
masks overlap substantially (any two masks share roughly
$1 - 2 \times 0.3 + 0.3^2 = 0.49$ of their active neurons at $p = 0.3$).
With $K = 6$ TTA augmentations, each is genuinely different (identity, flip,
rotate), yielding higher information content per sample.

\subsection{Why adaptive deferral outperforms global thresholding}

Global thresholding applies a single $\tau$ across all images. In a dataset
with variable image difficulty, this creates two failure modes: on easy images,
few pixels exceed $\tau$ and deferral contributes little; on hard images, many
pixels exceed $\tau$ and the system defers too aggressively, wasting review
capacity on pixels that are uncertain but correct.

Adaptive thresholding normalizes by image, deferring the top $\alpha$\%
most uncertain pixels regardless of absolute uncertainty magnitude. This
guarantees proportional contribution from each image and requires no
ground truth at test time.

The confidence-aware strategy addresses a different failure mode. Both global
and adaptive thresholding defer based on uncertainty alone, but uncertainty
and error are not identical. A pixel can be uncertain (moderate variance
across passes) yet correctly predicted (mean probability far from 0.5). The
confidence-aware score $s = u \cdot (1 - c)$ suppresses deferral for such
pixels, concentrating the budget on pixels that are both uncertain and
borderline. In our results, this produces the best low-budget efficiency
(ERR/Def.\% of 4.56 for confidence-aware vs.\ 4.19 for global TTA).

The adaptive and confidence-aware strategies are complementary. Adaptive
deferral works best at high deferral rates (25--30\%), where per-image
normalization prevents easy images from being ignored. Confidence-aware
deferral works best at low deferral rates (5--12\%), where concentrating
on the highest-risk pixels matters most. A hybrid policy---adaptive
thresholding with confidence weighting---is a natural extension we leave
for future work.

\subsection{Why calibration fails to improve deferral}

The mechanism is explained in Section~\ref{sec:calibration}: temperature
scaling is a monotonic transform that preserves uncertainty rankings. But
there is a deeper issue. ECE and deferral quality optimize genuinely
different objectives. ECE asks that errors be distributed proportionally
across confidence bins. Deferral asks that errors concentrate at high
uncertainty. A model that satisfies both simultaneously would need errors
to be both uniformly spread (for ECE) and concentrated (for deferral)---a
contradiction except in trivial cases.

Laves et al.\ \citep{laves2020well} observed this tension for classification.
Our results confirm it holds for dense pixel-level prediction. The practical
takeaway: report calibration and decision metrics side by side. Unc-AUROC,
risk-coverage curves, and achieved operating points answer a different
question from ECE.

\subsection{Deployment recommendations}

Based on the results in Sections~\ref{sec:results}--\ref{sec:ablations},
a practical deployment should:

\begin{enumerate}[leftmargin=*]
\item Use TTA (6 geometric transforms) for uncertainty when the goal is
pixel-level review triage. It dominates MC Dropout on both quality and speed
and requires no architectural changes.

\item Use confidence-aware deferral for review budgets $\leq$ 15\% of pixels,
adaptive thresholding for larger budgets. Both are computed from the existing
prediction and uncertainty maps at negligible cost.

\item Skip temperature scaling for deferral. If calibrated probabilities are
needed for other purposes (e.g., reporting confidence), compute them
separately; the deferral policy should use uncalibrated uncertainty.

\item Monitor Unc-AUROC on a small labeled subset when deploying across sites.
If it drops below $\sim$0.70, switch from pixel-level deferral to image-level
flagging.

\item Match the deferral rate to the workflow: around 10--12\% if review
budget is tight, around 25\% if maximizing error removal matters more than
coverage.
\end{enumerate}
